% !TEX program = lualatex
\documentclass[professionalfonts]{beamer}
\newcommand{\slidemode}{1}  % カラー投影モード
\usepackage{beamer_template}
\usepackage{enumerate}

\newcommand{\LectureNum}{14}
\newcommand{\LectureTitle}{スペクトル}
\newcommand{\CourseName}{応用数学}
\renewcommand{\TermName}{前期}

\begin{document}

\begin{frame}{第14回　スペクトル}
  \begin{block}{本回の内容}
  スペクトル
  \end{block}
  \vfill\hfill{\scriptsize \textcolor{gray}{第3章 フーリエ解析　§2}}
\end{frame}

\begin{frame}{定義：線スペクトル}
  \begin{block}{}
  周期 $2l$ の関数 $f(x) = \sum_{n=-\infty}^\infty c_{n} e^{in\pi x/l}$ に対し、\\[4pt]
  角周波数 $\omega_{n} = n\pi/l$ （とびとびの値）における $|c_{n}|$ を\\[4pt]
  \textbf{線スペクトル}（振幅スペクトル）という
  \end{block}
  {\footnotesize \textcolor{gray}{周期関数 $\to$ 離散的な $\omega_{n}$ に対応する離散スペクトル}}
\end{frame}

\begin{frame}{定義：連続スペクトル}
  \begin{block}{}
  非周期関数 $f(x)$ のフーリエ変換を $F(\omega)$ とするとき\\[4pt]
  \[
    S(\omega) = |F(\omega)| / \pi
  \]
  を \textbf{連続スペクトル}（振幅スペクトル）という
  \end{block}
  {\footnotesize \textcolor{gray}{周期 $2l$ の関数の線スペクトル $|c_{n}|$ は $\omega_{n}=n\pi/l$ 上で定義され、間隔 $\Delta\omega = \pi/l$。}\\
  \textcolor{gray}{$l\to\infty$ で $2l\, c_{n} \to F(\omega)$ より $|c_{n}|/\Delta\omega \to |F(\omega)|/(2\pi)$。}\\
  \textcolor{gray}{線スペクトルは正負両方の $n$ に立つので、片側に集約して $|F(\omega)|/\pi$ とする。}}
\end{frame}

\begin{frame}{例2　（p.\,100）}
  \begin{exampleblock}{問題}
    周期4の矩形波のスペクトルを求めよ\\[2pt]
    $f(x) = \begin{cases} 1 & (-1 < x < 1) \\ 1/2 & (x = 1, 3) \\ 0 & (1 < x < 3) \end{cases},  f(x+4) = f(x)$\\[2pt]
  \end{exampleblock}
\end{frame}

\begin{frame}{例2　解答}
  $l=2$ より $c_{n} = (1/4)\int_{-2}^{2} f(x) e^{-in\pi x/2} dx = (1/4)\int_{-1}^{1} e^{-in\pi x/2} dx$\\[4pt]
  $n=0: c_{0} = (1/4)\int_{-1}^{1} dx = 1/2$\\[4pt]
  $n\neq 0: c_{n} = (1/4)[-2 e^{-in\pi x/2}/(in\pi)]_{-1}^{1}$\\[4pt]
  $\hphantom{n\neq 0: c_{n}} = (-1/(2in\pi))(e^{-in\pi/2} - e^{in\pi/2}) = \sin(n\pi/2)/(n\pi)$\\[4pt]
  $n$ 奇数 $: |c_{n}| = 1/(|n|\pi),\quad n$ 偶数 ($\neq 0$) $: c_{n} = 0$
  \\[8pt]\textcolor{red}{\textbf{答}}：線スペクトル\\
  $|c_{0}| = 1/2,\ |c_{n}| = 1/(|n|\pi)$ （ $n=\pm 1,\pm 3,\ldots$ ）, $0$ （ $n=\pm 2,\pm 4,\ldots$ ）
\end{frame}

\begin{frame}{例3　（p.\,101）}
  \begin{exampleblock}{問題}
    次の関数のスペクトルを求めよ\\[2pt]
    $f(x) = \begin{cases} 1 & (|x| < 1) \\ 1/2 & (|x| = 1) \\ 0 & (|x| > 1) \end{cases}$\\[2pt]
  \end{exampleblock}
\end{frame}

\begin{frame}{例3　解答}
  非周期関数 $\to$ フーリエ変換 $F(u) = \int_{-\infty}^{\infty} f(x) e^{-iux} dx$\\[4pt]
  $F(u) = \int_{-1}^{1} e^{-iux} dx = [-e^{-iux}/(iu)]_{-1}^{1}$\\[4pt]
  $\hphantom{F(u)} = (e^{iu} - e^{-iu})/(iu) = 2\sin u / u$\\[4pt]
  $S(u) = |F(u)|/\pi = 2|\sin u|/(\pi |u|)$
  \\[8pt]\textcolor{red}{\textbf{答}}：連続スペクトル $S(u) = 2|\sin u|/(\pi |u|)$
\end{frame}

\begin{frame}{演習問題}
  \begin{exampleblock}{自分で解いてみよう}
  \begin{description}
    \item[問10] 次の関数のスペクトルを求めよ
    $f(x) = 1 - |x|  (|x| \leq 1),  f(x+2) = f(x)$\\
    \item[問11] 次の関数のスペクトルを求めよ
    $f(x) = \begin{cases} 1-|x| & (|x| \leq 1) \\ 0 & (|x| > 1) \end{cases}$\\
  \end{description}
  \end{exampleblock}
\end{frame}

\begin{frame}{問10　解答}
  周期 $T=2, l=1$ の偶関数 $: b_{n}=0$\\[4pt]
  $c_{0} = \int_{0}^1(1-x)dx = 1/2$\\[4pt]
  $a_{n} = 2\int_{0}^1(1-x)\cos(n\pi x)dx = 2(1-(-1)^n)/(n^{2}\pi^{2})$\\[4pt]
  n が奇数 $: a_{n} = 4/(n^{2}\pi^{2}), n$ が偶数 $: a_{n} = 0$\\[4pt]
  複素フーリエ係数 $: |c_{n}| = a_{n}/2 = 2/(n^{2}\pi^{2})$ （奇数 n ）
  \\[8pt]\textcolor{red}{\textbf{答}}：スペクトル $: |c_{0}|=1/2, |c_{n}|=2/(n^{2}\pi^{2})$ （ $n=1,3,5,...$ ） , 0 （ $n=2,4,6,...$ ）
\end{frame}

\begin{frame}{問11　解答}
  非周期関数 $ \to $ フーリエ変換でスペクトルを求める\\[4pt]
  $f(x)$ は偶関数 $: F(u) = 2\int_{0}^1(1-x)\cos(ux)dx$\\[4pt]
  $= 2\left\{[(1-x)\sin(ux)/u]_{0}^1 + (1/u)\int_{0}^1 \sin(ux)dx\right\}$\\[4pt]
  $= 2(1/u)[-\cos(ux)/u]_{0}^1 = 2(1-\cos u)/u^{2} = 4\sin^{2}(u/2)/u^{2}$
  \\[8pt]\textcolor{red}{\textbf{答}}：スペクトル $S(u) = |F(u)|/\pi = 4\sin^{2}(u/2)/(\pi u^{2})$ （連続スペクトル）
\end{frame}

\end{document}
